
\documentclass{article}
\usepackage{colortbl}
\usepackage{makecell}
\usepackage{multirow}
\usepackage{supertabular}

\begin{document}

\newcounter{utterance}

\centering \large Interaction Transcript for game `dond', experiment `semi\_en', episode 5 with Qwen3{-}32B{-}t0.0.
\vspace{24pt}

{ \footnotesize  \setcounter{utterance}{1}
\setlength{\tabcolsep}{0pt}
\begin{supertabular}{c@{$\;$}|p{.15\linewidth}@{}p{.15\linewidth}p{.15\linewidth}p{.15\linewidth}p{.15\linewidth}p{.15\linewidth}}
    \# & $\;$A & \multicolumn{4}{c}{Game Master} & $\;\:$B\\
    \hline

    \theutterance \stepcounter{utterance}  
    & & \multicolumn{4}{p{0.6\linewidth}}{
        \cellcolor[rgb]{0.9,0.9,0.9}{
            \makecell[{{p{\linewidth}}}]{
                \texttt{\tiny{[P1$\langle$GM]}}
                \texttt{You are playing a negotiation game in which you have to agree on how to divide a set of items among you and another player.} \\
\\ 
\texttt{Rules:} \\
\texttt{(a) You and the other player are given a set of items. Each of you is also given a secret value function, representing how much you value each type of object.} \\
\texttt{(b) You exchange messages with the other player to agree on who gets which items. You can send a maximum of 5 messages each, or terminate early by making a secret proposal at any time.} \\
\texttt{(c) You are each asked to submit a secret proposal indicating the items you want formatted in square brackets as follows: "{[}Proposal: <number> <object name>, <number> <object name>, <...>{]}"} \\
\texttt{(d) If your proposals are complementary, i.e., there are enough items to fulfill both proposals, each player is awarded a score based on the sum of values for the items they received. Otherwise, both of you get zero points.} \\
\texttt{(e) Your goal is to maximize the score you receive.} \\
\\ 
\texttt{Let us start.} \\
\\ 
\texttt{The set of available items is:} \\
\\ 
\texttt{2 dice, 1 ball, 1 glue, 1 battery} \\
\\ 
\texttt{Your secret value function is:} \\
\\ 
\texttt{dice: 1, ball: 7, glue: 0, battery: 1} \\
\\ 
\texttt{IMPORTANT: Your messages, unless it is the secret proposal, are directly transmitted to the other player, so do not include any response to the rules or text announcing your message. To make a secret proposal, use the indicated format. Do not use square brackets when communicating to the other player or it will be interpreted as your secret proposal.} \\
            }
        }
    }
    & \\ \\

    \theutterance \stepcounter{utterance}  
    & \multicolumn{4}{p{0.6\linewidth}}{
        \cellcolor[rgb]{0.9,0.9,0.9}{
            \makecell[{{p{\linewidth}}}]{
                \texttt{\tiny{[P1$\rangle$GM]}}
                \texttt{<think>} \\
\texttt{Okay, let's see. I need to negotiate the division of these items: 2 dice, 1 ball, 1 glue, 1 battery. My values are dice:1, ball:7, glue:0, battery:1. So the ball is super valuable to me, while glue is worthless. The other player's values are secret, so I don't know what they want. } \\
\\ 
\texttt{First, I should figure out what the other player might value. Since I don't have their value function, I need to make some assumptions. Maybe they value the glue or battery more? Or maybe they also want the ball? But since the ball is only one, if both of us want it, we can't both get it. } \\
\\ 
\texttt{My main goal is to get the ball because it's worth 7 points, which is way more than the others. The dice and battery are each worth 1, so maybe I can get one of those if needed. But I need to make sure that my proposal doesn't conflict with the other player's. If I take the ball, maybe they can take the other items. Let me think about possible proposals.} \\
\\ 
\texttt{If I propose taking the ball, then the other player would need to take the remaining items: 2 dice, 1 glue, 1 battery. But if their value function is such that they value, say, glue or battery highly, maybe that's okay. But if they also want the ball, then we can't both have it. So maybe I should try to negotiate by offering something in exchange for the ball, but since I don't know their values, it's tricky.} \\
\\ 
\texttt{Alternatively, maybe I can propose taking the ball and one die, and let the other player take the other die, glue, and battery. But that might not be complementary. Wait, the total items are 2 dice, 1 ball, 1 glue, 1 battery. If I take the ball and one die, the other player would have one die, glue, battery. That's possible. But if the other player's proposal is to take, say, the ball, then our proposals would conflict. So maybe I should make a proposal that's more likely to be complementary.} \\
\\ 
\texttt{Another approach: since the ball is unique, if I take it, the other player can't. So maybe I should propose taking the ball and then some low{-}value items for me, but maybe high for them. For} \\
            }
        }
    }
    & & \\ \\

    \theutterance \stepcounter{utterance}  
    & & & \multicolumn{2}{p{0.3\linewidth}}{
        \cellcolor[rgb]{0.9,0.9,0.9}{
            \makecell[{{p{\linewidth}}}]{
                \texttt{\tiny{[GM$|$GM]}}
                \texttt{continue} \\
            }
        }
    }
    & & \\ \\

    \theutterance \stepcounter{utterance}  
    & & & \multicolumn{2}{p{0.3\linewidth}}{
        \cellcolor[rgb]{0.9,0.9,0.9}{
            \makecell[{{p{\linewidth}}}]{
                \texttt{\tiny{[GM$|$GM]}}
                \texttt{<think>} \\
\texttt{Okay, let's see. I need to negotiate the division of these items: 2 dice, 1 ball, 1 glue, 1 battery. My values are dice:1, ball:7, glue:0, battery:1. So the ball is super valuable to me, while glue is worthless. The other player's values are secret, so I don't know what they want. } \\
\\ 
\texttt{First, I should figure out what the other player might value. Since I don't have their value function, I need to make some assumptions. Maybe they value the glue or battery more? Or maybe they also want the ball? But since the ball is only one, if both of us want it, we can't both get it. } \\
\\ 
\texttt{My main goal is to get the ball because it's worth 7 points, which is way more than the others. The dice and battery are each worth 1, so maybe I can get one of those if needed. But I need to make sure that my proposal doesn't conflict with the other player's. If I take the ball, maybe they can take the other items. Let me think about possible proposals.} \\
\\ 
\texttt{If I propose taking the ball, then the other player would need to take the remaining items: 2 dice, 1 glue, 1 battery. But if their value function is such that they value, say, glue or battery highly, maybe that's okay. But if they also want the ball, then we can't both have it. So maybe I should try to negotiate by offering something in exchange for the ball, but since I don't know their values, it's tricky.} \\
\\ 
\texttt{Alternatively, maybe I can propose taking the ball and one die, and let the other player take the other die, glue, and battery. But that might not be complementary. Wait, the total items are 2 dice, 1 ball, 1 glue, 1 battery. If I take the ball and one die, the other player would have one die, glue, battery. That's possible. But if the other player's proposal is to take, say, the ball, then our proposals would conflict. So maybe I should make a proposal that's more likely to be complementary.} \\
\\ 
\texttt{Another approach: since the ball is unique, if I take it, the other player can't. So maybe I should propose taking the ball and then some low{-}value items for me, but maybe high for them. For} \\
            }
        }
    }
    & & \\ \\

    \theutterance \stepcounter{utterance}  
    & & \multicolumn{4}{p{0.6\linewidth}}{
        \cellcolor[rgb]{0.9,0.9,0.9}{
            \makecell[{{p{\linewidth}}}]{
                \texttt{\tiny{[GM$\rangle$P2]}}
                \texttt{You are playing a negotiation game in which you have to agree on how to divide a set of items among you and another player.} \\
\\ 
\texttt{Rules:} \\
\texttt{(a) You and the other player are given a set of items. Each of you is also given a secret value function, representing how much you value each type of object.} \\
\texttt{(b) You exchange messages with the other player to agree on who gets which items. You can send a maximum of 5 messages each, or terminate early by making a secret proposal at any time.} \\
\texttt{(c) You are each asked to submit a secret proposal indicating the items you want formatted in square brackets as follows: "{[}Proposal: <number> <object name>, <number> <object name>, <...>{]}"} \\
\texttt{(d) If your proposals are complementary, i.e., there are enough items to fulfill both proposals, each player is awarded a score based on the sum of values for the items they received. Otherwise, both of you get zero points.} \\
\texttt{(e) Your goal is to maximize the score you receive.} \\
\\ 
\texttt{Let us start.} \\
\\ 
\texttt{The set of available items is:} \\
\\ 
\texttt{2 dice, 1 ball, 1 glue, 1 battery} \\
\\ 
\texttt{Your secret value function is:} \\
\\ 
\texttt{dice: 0, ball: 2, glue: 8, battery: 0} \\
\\ 
\texttt{IMPORTANT: Your messages, unless it is the secret proposal, are directly transmitted to the other player, so do not include any response to the rules or text announcing your message. To make a secret proposal, use the indicated format. Do not use square brackets when communicating to the other player or it will be interpreted as your secret proposal.} \\
\\ 
\\ 
\texttt{<think>} \\
\texttt{Okay, let's see. I need to negotiate the division of these items: 2 dice, 1 ball, 1 glue, 1 battery. My values are dice:1, ball:7, glue:0, battery:1. So the ball is super valuable to me, while glue is worthless. The other player's values are secret, so I don't know what they want. } \\
\\ 
\texttt{First, I should figure out what the other player might value. Since I don't have their value function, I need to make some assumptions. Maybe they value the glue or battery more? Or maybe they also want the ball? But since the ball is only one, if both of us want it, we can't both get it. } \\
\\ 
\texttt{My main goal is to get the ball because it's worth 7 points, which is way more than the others. The dice and battery are each worth 1, so maybe I can get one of those if needed. But I need to make sure that my proposal doesn't conflict with the other player's. If I take the ball, maybe they can take the other items. Let me think about possible proposals.} \\
\\ 
\texttt{If I propose taking the ball, then the other player would need to take the remaining items: 2 dice, 1 glue, 1 battery. But if their value function is such that they value, say, glue or battery highly, maybe that's okay. But if they also want the ball, then we can't both have it. So maybe I should try to negotiate by offering something in exchange for the ball, but since I don't know their values, it's tricky.} \\
\\ 
\texttt{Alternatively, maybe I can propose taking the ball and one die, and let the other player take the other die, glue, and battery. But that might not be complementary. Wait, the total items are 2 dice, 1 ball, 1 glue, 1 battery. If I take the ball and one die, the other player would have one die, glue, battery. That's possible. But if the other player's proposal is to take, say, the ball, then our proposals would conflict. So maybe I should make a proposal that's more likely to be complementary.} \\
\\ 
\texttt{Another approach: since the ball is unique, if I take it, the other player can't. So maybe I should propose taking the ball and then some low{-}value items for me, but maybe high for them. For} \\
            }
        }
    }
    & \\ \\

    \theutterance \stepcounter{utterance}  
    & & & \multicolumn{4}{p{0.6\linewidth}}{
        \cellcolor[rgb]{0.9,0.9,0.9}{
            \makecell[{{p{\linewidth}}}]{
                \texttt{\tiny{[GM$\langle$P2]}}
                \texttt{<think>} \\
\texttt{Okay, let's start by understanding the problem. I need to divide the items between me and another player. The items are 2 dice, 1 ball, 1 glue, and 1 battery. My value function is dice:1, ball:7, glue:0, battery:1. The other player's values are secret. The key is to make a proposal that's complementary to theirs so we both get points. If not, we both get zero.} \\
\\ 
\texttt{First, my highest value item is the ball. So ideally, I want the ball. But if the other player also wants the ball, our proposals would conflict. So maybe I should try to get the ball and then let the other player take the rest. But I need to consider what the other player might value. Since I don't know their values, I have to make educated guesses.} \\
\\ 
\texttt{The other player might value glue or battery more. For example, if they value glue highly, they might want that. If they value battery, same. The dice are low value for me, so maybe I can offer them to the other player. Let me think of possible proposals.} \\
\\ 
\texttt{If I propose taking the ball and one die, then the other player could take the other die, glue, and battery. That would leave all items distributed. But if the other player's proposal is to take the ball, then we conflict. So maybe I should propose taking the ball and one die, and the other player takes the rest. Alternatively, maybe I take the ball and battery, and the other player takes the dice and glue. But battery is low value for me, so maybe that's better.} \\
\\ 
\texttt{Wait, my battery is worth 1, same as a die. So taking the ball (7) plus battery (1) gives me 8. If the other player takes the two dice and glue, maybe that's acceptable. But if the other player's values are such that they need the battery or something else, maybe that's okay. } \\
\\ 
\texttt{Alternatively, maybe I take the ball and one die, and the other player takes the other die, glue, and battery. That way, I have 7+1=8, and the other player has 1 die (maybe they value it?), glue, and battery. If their values are such that they can get points from those, then our proposals are complementary.} \\
\\ 
\texttt{But how do I communicate this? Since I can send up to 5 messages,} \\
            }
        }
    }
     \\ \\

    \theutterance \stepcounter{utterance}  
    & & & \multicolumn{2}{p{0.3\linewidth}}{
        \cellcolor[rgb]{0.9,0.9,0.9}{
            \makecell[{{p{\linewidth}}}]{
                \texttt{\tiny{[GM$|$GM]}}
                \texttt{continue} \\
            }
        }
    }
    & & \\ \\

    \theutterance \stepcounter{utterance}  
    & & & \multicolumn{2}{p{0.3\linewidth}}{
        \cellcolor[rgb]{0.9,0.9,0.9}{
            \makecell[{{p{\linewidth}}}]{
                \texttt{\tiny{[GM$|$GM]}}
                \texttt{<think>} \\
\texttt{Okay, let's start by understanding the problem. I need to divide the items between me and another player. The items are 2 dice, 1 ball, 1 glue, and 1 battery. My value function is dice:1, ball:7, glue:0, battery:1. The other player's values are secret. The key is to make a proposal that's complementary to theirs so we both get points. If not, we both get zero.} \\
\\ 
\texttt{First, my highest value item is the ball. So ideally, I want the ball. But if the other player also wants the ball, our proposals would conflict. So maybe I should try to get the ball and then let the other player take the rest. But I need to consider what the other player might value. Since I don't know their values, I have to make educated guesses.} \\
\\ 
\texttt{The other player might value glue or battery more. For example, if they value glue highly, they might want that. If they value battery, same. The dice are low value for me, so maybe I can offer them to the other player. Let me think of possible proposals.} \\
\\ 
\texttt{If I propose taking the ball and one die, then the other player could take the other die, glue, and battery. That would leave all items distributed. But if the other player's proposal is to take the ball, then we conflict. So maybe I should propose taking the ball and one die, and the other player takes the rest. Alternatively, maybe I take the ball and battery, and the other player takes the dice and glue. But battery is low value for me, so maybe that's better.} \\
\\ 
\texttt{Wait, my battery is worth 1, same as a die. So taking the ball (7) plus battery (1) gives me 8. If the other player takes the two dice and glue, maybe that's acceptable. But if the other player's values are such that they need the battery or something else, maybe that's okay. } \\
\\ 
\texttt{Alternatively, maybe I take the ball and one die, and the other player takes the other die, glue, and battery. That way, I have 7+1=8, and the other player has 1 die (maybe they value it?), glue, and battery. If their values are such that they can get points from those, then our proposals are complementary.} \\
\\ 
\texttt{But how do I communicate this? Since I can send up to 5 messages,} \\
            }
        }
    }
    & & \\ \\

    \theutterance \stepcounter{utterance}  
    & & \multicolumn{4}{p{0.6\linewidth}}{
        \cellcolor[rgb]{0.9,0.9,0.9}{
            \makecell[{{p{\linewidth}}}]{
                \texttt{\tiny{[P1$\langle$GM]}}
                \texttt{<think>} \\
\texttt{Okay, let's start by understanding the problem. I need to divide the items between me and another player. The items are 2 dice, 1 ball, 1 glue, and 1 battery. My value function is dice:1, ball:7, glue:0, battery:1. The other player's values are secret. The key is to make a proposal that's complementary to theirs so we both get points. If not, we both get zero.} \\
\\ 
\texttt{First, my highest value item is the ball. So ideally, I want the ball. But if the other player also wants the ball, our proposals would conflict. So maybe I should try to get the ball and then let the other player take the rest. But I need to consider what the other player might value. Since I don't know their values, I have to make educated guesses.} \\
\\ 
\texttt{The other player might value glue or battery more. For example, if they value glue highly, they might want that. If they value battery, same. The dice are low value for me, so maybe I can offer them to the other player. Let me think of possible proposals.} \\
\\ 
\texttt{If I propose taking the ball and one die, then the other player could take the other die, glue, and battery. That would leave all items distributed. But if the other player's proposal is to take the ball, then we conflict. So maybe I should propose taking the ball and one die, and the other player takes the rest. Alternatively, maybe I take the ball and battery, and the other player takes the dice and glue. But battery is low value for me, so maybe that's better.} \\
\\ 
\texttt{Wait, my battery is worth 1, same as a die. So taking the ball (7) plus battery (1) gives me 8. If the other player takes the two dice and glue, maybe that's acceptable. But if the other player's values are such that they need the battery or something else, maybe that's okay. } \\
\\ 
\texttt{Alternatively, maybe I take the ball and one die, and the other player takes the other die, glue, and battery. That way, I have 7+1=8, and the other player has 1 die (maybe they value it?), glue, and battery. If their values are such that they can get points from those, then our proposals are complementary.} \\
\\ 
\texttt{But how do I communicate this? Since I can send up to 5 messages,} \\
            }
        }
    }
    & \\ \\

    \theutterance \stepcounter{utterance}  
    & \multicolumn{4}{p{0.6\linewidth}}{
        \cellcolor[rgb]{0.9,0.9,0.9}{
            \makecell[{{p{\linewidth}}}]{
                \texttt{\tiny{[P1$\rangle$GM]}}
                \texttt{<think>} \\
\texttt{Okay, let's try to figure out the best way to approach this negotiation. My main goal is to get the ball because it's worth 7 points, which is way more than the other items. The other player's values are secret, so I don't know what they want. But I need to make sure that my proposal doesn't conflict with theirs. If we both want the same item, we both get zero. So I need to find a way to get the ball without the other player also wanting it.} \\
\\ 
\texttt{First, maybe I should propose taking the ball and one of the low{-}value items for me, like a die or the battery. That way, I can get the ball and leave the other items for the other player. Let's see: if I take the ball and one die, then the remaining items are one die, glue, and battery. The other player could take those. But if the other player's value function is such that they need the ball, then we'd conflict. But since I don't know their values, I have to assume they might not want the ball. Alternatively, maybe they value the glue or battery more.} \\
\\ 
\texttt{Another thought: the battery is worth 1 to me, same as a die. So maybe I can take the ball and the battery, leaving the two dice and glue for the other player. That way, I get 7+1=8 points. The other player might get value from the dice or glue. If their values are such that they can get points from those, then our proposals are complementary.} \\
\\ 
\texttt{But how do I communicate this? Since I can send up to 5 messages, maybe I should start by proposing a division where I take the ball and one other item, and the other player takes the rest. Let me try that. So my first message could be: "I would like the ball and one die. You can have the other die, glue, and battery." But I need to format the secret proposal correctly. Wait, the instructions say that secret proposals should be in the format {[}Proposal: ...{]}. So maybe I should make a secret proposal right away. Let me check the rules again.} \\
\\ 
\texttt{Rule (c) says that secret proposals are submitted in the format {[}Proposal: ...{]}. So if I make a secret proposal, it's a final offer. But before that, I can exchange messages. Maybe I should first send a message to the other player to suggest a division, then} \\
            }
        }
    }
    & & \\ \\

    \theutterance \stepcounter{utterance}  
    & & & \multicolumn{2}{p{0.3\linewidth}}{
        \cellcolor[rgb]{0.9,0.9,0.9}{
            \makecell[{{p{\linewidth}}}]{
                \texttt{\tiny{[GM$|$GM]}}
                \texttt{proposal does not include number/name pairs} \\
            }
        }
    }
    & & \\ \\

\end{supertabular}
}

\end{document}
